\documentclass[12pt]{article}

\usepackage[a4paper, total={6in, 8in}]{geometry}
\usepackage{textcomp}

\begin{document}
\setlength{\parindent}{0pt}

\def \t {\textrightarrow}
\def \v {\vspace{1em}}
\def \bi {\begin{itemize}}
\def \ei {\end{itemize}}
\def \s[#1] {\section*{#1}}
\def \ss[#1] {\subsection*{#1}}
\def \sss[#1] {\subsubsection*{#1}}

\s[La sera fiesolana]
Appunti di oggi presi un po male, chiedere \\
La natura come donna dirompente, non come mulier \\
Opera di ispirazione francescana \t fa parte di Alcyone e ha scansione che ricorda il cantico delle natura

\v

Incipit presenta un iperbato \t ed enjambemnt fra primo e secondo verso \t parole fresche è sinestesia \\
La similitudine chiarifica la lontanza delle immagini \\
Analogia e onomatopea sono le figure retoriche fondamentali di Pascoli \\
Immagie rarefatta dell'oggettistica di questa sera a fiesole \\
Il fruscio del gelso nelle mani di chi lo coglie \t lui osserva figura non identificata \t gelso produce le more di gelso \t con onomatopea evidente \\
La scala è alta, che s'annera contro il fusto che s'inargenta \t sembra risultare scura in alto, mentre la luce argentea che coglie il fusto riceve dalla luna \\
Il suo linguaggio non parla di simboli, come nello stile decadentista \t lo dice Verga, e quindi dimostra che in Sperelli c'è l'inetto \\
L'alone che contorna l'astro in cielo si vede quando la notte è serena \\
Versi 11 \t anafora \\
La campagna che sente = emblema di una umanizzazione della natura \t luna è sempre confortante, e spettatrice ma non interlocutrici, e anche adutrix = aiutante

\v

Versi 29-30-31 sono fortemente allusivi \t si vede poi perche \\
Laudata si \t san francesco \\
Viso di perla riferito alla sera = umanizzazione della natura \t analogia, e anche sinestesia \\
La sera è una donna che provoca attraverso le sue sembianze \\
Questi umidi occhi \t perche dovrebbero essere umidi? perche emozionati e in preda a sentimento dell'animo \\
Dolci le mie parole \t ripetizione nei versi precedenti, sinestesia e anastrofe \\
Il tono è conativo/imperativo/esortativo? \t si è anche un po conativo \\
Bruiva è onomatopea \t bruiva ha significato quasi opposto da scaldare \t contrapposta alla pioggia fredda? \\
Commiato è termine molto elettivo \\
Primvera che sfocia nell'estate \t con un commiato lacrimoso, come se non volesse andarsene \t è la stagione della vita \\
I pini novelli \t utilizza un suo linguaggio \\
Palazzeschi ribalta il dialogo/monolgo tra Ermione e d'Annunzio

\v

Alcune immagini rimandono a un po leopardi \t come la mano che tiene il gelso, che probabilmente lavora \t l'operosita abruzzese viene colta con uno sguardo novecentesco e rimanda a leopardi \\
IMP. ESAME \t anche in Lavandare emigrazione e operosita, e rimanda ai pescatori di Verga e al mondo di Nedda \\
Natura uomo, sensualismo e immagine candida e pura sono al centro \\
C'è una enumeratio e climax ascendente in cui elenca gelsi, olmi etc. \\
Decadentismo presente nel grano non biondo e non verde \t è un indefinito, riminescenza leopardiana \t la natura vista in tutte le sue sfumature \\
Il fieno \t descrive la violenza dell'agricoltore che usa la falce per tagliare il grano \t e in falce di luna calante parla di luna che compare a spicchi \\
Il fieno patisce \t violenza dell'uomo sulla natura \\
Messi in risalto gli ulivi \t pianta legata alla sacralita di cristo, ma è anche una pianta dal tronco nodoso e non regolare \t quindi difficile da vedere in modo rettilineo, ed evoca quindi le immagini di un corpo umano \\
L'uomo vecchio si puo essere paragonato alla nodosita dell ulivo \t qua si presenta in modo blasfemo perche? \\
La vita monastica è anche ombra \t come nel caso della monaca di monza, per cui la vita in monastero puo essere vivace \\
È presente quindi una piccante insinuazione riguardo alla pedofilia nella chiesa ?

\v

Per il cinto che ti cinge \t i lacci che circondano i fasci di fieno sono paragonati al cinto sotto il seno della donna che è la sera \\
Amar d'amore è un gioco semantico \\
Segreto sacro dei monti \t la natura ha dei segreti che non rivela all uomo \\
Il filtro delle labbra paragonato al profilo delle colline \\
Cambio tonale che riporta alla virginita del canto francescano negli ultimi versi \t palpitare delle stelle in attesa della notte fa vedere natura umanizzata

\v

Testo mischia sacro e profano, puro e voluttuoso
\end{document}
